\chapter{生命系统的多尺度适应动力学：从局部可塑性到种群演化}
\label{chp:life_manifold_dynamics}

生命系统同时包含毫秒级状态传播、分钟至日尺度的可塑性、跨睡眠周期的记忆巩固，以及跨世代的种群变化。把这些过程分别写成神经电位方程、学习规则、网络重构模型和群体遗传方程，是必要的专业分工；但若因此把各尺度视为互不相干的孤岛，我们就无法说明快速活动如何留下慢结构。相反，若把它们宣称为同一个物理场方程的严格同构，又会抹去对象类型、时间单位、因果通道和证据等级的差异。本章把生命适应重建为一组通过统计量、约束与读出接口耦合的计算动力学。HSF-HD提供一般状态动力学的组织语言，而不替代神经科学或演化生物学；AgentOS的快态、情景记录与慢结构只作为工程实例。

\section{跨尺度模型族：对象、时钟与接口}
\label{sec:breaking_scale_barriers}
\label{sec:unified_action_across_scales}

我们定义尺度集合 $K=\{n,s,c,p\}$，分别代表神经活动、突触结构、系统级连接和种群分布。每一尺度拥有自己的状态空间与时钟：神经活动 $x_n(t)\in\mathbb{R}^{d_n}$ 以秒为单位；突触参数 $W_s(\tau_s)\in\mathcal W$ 以分钟至日为单位；系统级结构 $G_c(\tau_c)\in\mathcal G$ 以日月至生命周期为单位；种群分布 $p_p(\tau_p)\in\Delta^{m-1}$ 以世代为单位。$\mathcal W$ 可以是带符号权重矩阵或局部结构网络，$\mathcal G$ 则记录跨区域的有效连接、记忆索引和长期可达关系。四类对象不能因都可画成“流形”而互换；神经活动强度不是信息熵，突触参数不是焦耳能耗，种群频率也不是个体信念。

在这些类型约束下，多尺度系统写成耦合模型族
\begin{align}
 \dot x_n &= f_n(x_n,W_s,u_n,t)+\Sigma_n(x_n,t)\xi_n(t),\\
 \frac{dW_s}{d\tau_s} &= f_s(W_s;A_s[x_n],m_s,r_s),\\
 \frac{dG_c}{d\tau_c} &= f_c(G_c;A_c[W_s,x_n],z_c),\\
 \frac{dp_p}{d\tau_p} &= f_p(p_p;\bar q,G_e,\mu,N_e).
\end{align}
其中 $u_n$ 是外部输入，$m_s$ 是调制信号，$r_s$ 是局部资源状态，$z_c$ 是睡眠或系统门控，$\bar q$ 是表型适合度，$G_e$ 是遗传协方差或迁移结构，$\mu$ 是突变参数，$N_e$ 是有效种群大小。$A_s$ 与 $A_c$ 不是神秘的升维算子，而是有窗口、有损失的统计接口，例如脉冲相关、预测残差、回放次数和跨任务复现率。若窗口、采样率或身份绑定规则改变，它们的输出也会改变，这些变化必须进入模型。

为了比较而不混同各尺度，我们可定义有限时域目标 $\mathcal J=\sum_k\alpha_k\mathcal J_k+\sum_{k\to\ell}\beta_{k\ell}\mathcal C_{k\ell}$。$\mathcal J_k$ 是尺度内的预测、稳态、繁殖或资源代价，$\mathcal C_{k\ell}$ 是接口失配。系数只用于统一记账，不能据此推出自然界在求解一个真实的总作用量。若问题具备可微目标和确定边界，可以用变分或最优控制；若结构编辑离散、目标随环境变化，则应使用混合系统、随机过程或多目标决策。统一性是共同的状态—结构—环境闭环，不是方程形式严格同构。

反例表明接口不可省略：一次高频放电可能因门控关闭而不产生可塑性；一次突触增强可能在稳态调节中被抵消；个体学到的行为通常不会把获得性结构直接写入生殖系；种群适合度上升也不保证每个个体的神经网络更复杂。因此每条上行链必须给出来源、窗口、充分统计量和干预结果，每条下行链必须给出约束载体、作用延迟和替代解释。只有这样，多尺度动力学才是可检验模型，而不是把不同学科术语放进同一张比喻表。

局部结构网络与全局分布表示也在各尺度保持区分：前者保存对象身份、关系与可达路径，后者编码当前内容、统计相似性与不确定性。二者通过任务相关身份键、绑定算子、耦合窗口和读出接口联合。相同结构可承载不同状态，不同结构也可能产生相近读出；因此不能从脑区功能相似推出机制等同。HSF-HD关注接口条件，具体载体仍由领域模型和实验确定。

这种建模方式还要求为每个接口维护版本与不确定性。若低层统计来自不同个体、不同仪器或不同任务，必须先校准身份与量纲，再允许进入高层更新；若摘要的置信区间跨越提交阈值，高层只能保留候选而不能修改长期结构。上行统计失败时应回到原始事件或更细尺度复核，下行约束失败时应允许低层报告不可行原因。由此，多尺度系统不是单向层级命令，而是带回执、拒绝与修订的闭环协议。该协议既适用于生物测量，也适用于分布式人工智能系统，并能明确指出信息在何处丢失、责任在何处转移。

\section{微观层：受约束神经状态传播}
\label{sec:micro_neural_field_dynamics}

在毫秒尺度，我们研究膜电位、放电率或群体活动如何在既定连接上演化。设 $x_i(t)$ 是第 $i$ 个单元的无量纲活动，$W_{ij}$ 表示从 $j$ 到 $i$ 的有效耦合，$d_{ij}\ge0$ 是秒单位传输延迟，$u_i(t)$ 是经传感接口校准后的输入，$a_i(t)$ 是调质和行为状态。一个用于比较实验的随机延迟模型为
\begin{equation}
 \tau_i\dot x_i=-x_i+\phi_i\!\left(\sum_jW_{ij}x_j(t-d_{ij})+B_iu_i+M_ia_i-\theta_i\right)+\sigma_i\xi_i.
\end{equation}
其中 $\tau_i$ 是响应时间常数，$\phi_i$ 是有界增益，$\theta_i$ 是阈值，$\xi_i$ 是单位方差噪声。该式是计算模型，并非Hodgkin--Huxley方程的“量子化推广”。若问题涉及离子通道和动作电位波形，应采用带电导、反转电位和门控变量的生物物理模型；若只关心群体传播、选择和稳定性，上式才是合适的粗粒化。

固定 $W$ 与输入统计时，平衡点满足 $x^*=\phi(Wx^*+Bu+Ma-\theta)$。若 $\phi$ 的Lipschitz常数为 $L_\phi$ 且 $L_\phi\|W\|<1$，无延迟确定系统在相应范数下存在唯一收缩不动点；这只证明局部模型的状态收敛，不证明行为正确。延迟增大、增益升高或抑制减弱时，特征根可能越过虚轴而出现振荡甚至不稳定。所谓癫痫样失稳可以被操作化为谱半径、同步率和恢复时间的异常组合，但真实临床解释仍需电生理证据，不能由抽象稳定性条件直接诊断。

神经调质在模型中不是给状态注入目的的规范场，而是改变增益、阈值、可塑性资格和行动选择的可观测输入。多巴胺、去甲肾上腺素和血清素的作用依赖受体、区域与时间，既可能促进某类更新，也可能抑制另一类更新。我们把 $a_i(t)$ 当作条件变量，通过消融、药理干预或任务切换辨识 $M_i$，不把奖励、唤醒与价值混成一个标量。泄漏项的维度是活动每秒变化率，不是热力学熵产；实际焦耳能耗需由氧耗、葡萄糖代谢、电流或硬件功率单独测量。

结构矩阵 $W$ 给出谁能影响谁，分布活动 $x$ 编码当前内容，二者经节点身份、时间戳和读出矩阵 $R$ 结合，输出 $y=R(x,W,q)$，其中任务 $q$ 决定保留哪些差异。通过扰动 $W$、钳制 $x$ 与改变 $R$ 的对照实验，可以分别估计结构、状态与任务接口的贡献，而不从某个脑区类比推出工程机制等同。

离散仿真以步长 $\Delta t$ 更新 $x^{r+1}=x^r+\Delta t f_n(x^r,W,u^r)+\sqrt{\Delta t}\Sigma_n\epsilon^r$。步长须小于相关时间常数并位于积分器稳定域，延迟需要缓存和插值；否则数值爆炸不能解释为生物相变。微观层最终保留的是可跨载体比较的状态、结构、输入、调节、噪声与读出关系，同时明确其生物物理边界。

验证微观模型时，我们至少需要三组对照。第一组固定结构而改变输入和调质，检验状态传播与增益项；第二组固定输入统计而局部切断连接，检验网络路径和延迟；第三组改变读出任务但保持底层活动，检验所谓全局表征是否只是分析器选择的结果。评价指标包括一步预测、长时滚动误差、扰动恢复、校准度与跨个体复现，不能只报告训练拟合。若一个更简单的线性或广义线性模型已经解释数据，就不应凭复杂几何形式宣称更深机制。

微观状态还受边界与资源限制。传感饱和、代谢供应、抑制回路和身体动作都会改变可达状态集；环境反馈又通过新的观测进入下一轮更新。因而神经系统不是孤立求解器，而是具身闭环的一部分。HSF-HD在此只提出可观测状态、慢结构、输入、行动和边界之间的关系。具体系统是否采用脉冲、连续率、振荡或事件编码，必须由任务与证据决定，这也为后续讨论可塑性留下了明确接口。

\section{中观层：局部可塑性、资格迹与STDP}
\label{sec:meso_cefe_and_stdp}

在分钟至日尺度，快速活动不再是研究终点，而成为结构更新的证据。设 $W_{ij}$ 是有界突触效能，$e_{ij}(t)$ 是突触资格迹，记录局部前后活动是否在因果窗口内相遇；$m(t)$ 是延迟到达的调制或结果信号。我们定义
\begin{align}
 \tau_e\dot e_{ij}&=-e_{ij}+K_+(\Delta t_{ij})-K_-(\Delta t_{ji}),\\
 \frac{dW_{ij}}{d\tau_s}&=\eta_s m(\tau_s)e_{ij}-\lambda_s(W_{ij}-W_{ij}^{0})-\rho_sH_i(W)+\zeta_{ij}.
\end{align}
其中 $K_+$ 与 $K_-$ 是由实验拟合的时序窗口，$W^0$ 是参考连接，$H_i$ 表示稳态可塑性或资源竞争，$\zeta_{ij}$ 汇总未建模扰动。这个模型保留STDP的具体机制：前后脉冲时差影响资格，资格与结果调制相遇后才决定长期保留；它不把余弦相位差宣称为LTP与LTD的几何必然。

经典成对STDP曲线常用指数窗表示：突触前脉冲先于突触后脉冲时，$K_+(\Delta t)=A_+e^{-\Delta t/\tau_+}$；顺序相反时，$K_-(|\Delta t|)=A_-e^{-|\Delta t|/\tau_-}$。这一经验形式会受放电频率、脉冲三元组、树突位置、抑制背景、钙动力学和行为状态影响。“前早后晚必增强”只是特定制备与参数区间内的经验主张，不是普遍定理。若数据表现为对称Hebbian窗、反Hebbian窗或频率依赖翻转，模型应更换核函数或加入隐状态，而不能把反例解释成抽象流形的例外。

局部学习能否等同于目标函数梯度，需要额外条件。设行为风险为 $L(W)$，若资格迹是轨迹对参数得分函数的无偏估计，调制量是合适的回报预测误差，且采样分布覆盖相关轨迹，则三因子更新的期望可以近似策略梯度。这个数学结果不适用于所有STDP，也不意味着单个突触获得全局误差的精确偏导。反向传播依赖可微计算图、正向中间量和权重传递；局部可塑性依赖空间邻近、时序相关与调制广播。二者可在某些任务上产生相似下降方向，却不具有整体算法同构。

为防止无界强化，我们令 $W_{ij}\in[W^-_{ij},W^+_{ij}]$，离散更新后执行 $W^{r+1}=\Pi_{\mathcal W}(W^r+\Delta\tau_sF_s(W^r,e^r,m^r))$。当 $F_s$ 局部Lipschitz、步长足够小且投影集合闭凸时，迭代保持可行；这只保证参数不越界。学习是否改善任务，仍需留出数据、因果干预和长期保持测试。我们分别记录预测损失、连接稀疏度、更新次数和实际代谢或硬件功耗，避免把更新幅度变小误报为更节能。

生物实现还受突触饱和、蛋白合成、局部资源和稳态调节限制。同一相关模式若持续出现，可能先增强后因归一化而回落；一次强烈结果也可能因资格迹已经衰减而不被保留。实验必须同时扫描时差、结果延迟、重复频率和保持时长。只比较更新前后的平均权重会遗漏信用分配过程，也无法区分直接可塑性与网络重排造成的表观变化。

还要把“写入失败”与“读出失败”分开诊断。若刺激后连接指标改变而行为读出不变，可能是更新未进入任务相关通路；若行为短时改善但跨日消失，则更像稳定化不足；若替换结果信号后更新方向随之翻转，才为三因子信用分配提供较强证据。相反，若仅调整总体兴奋性也能复现实验效应，就不能把观察结果专归因于STDP。因而最小验证集应包含时序置乱、结果延迟置换、稳态阻断和身份错配四类对照，并在相同活动预算下比较保持与迁移。

在AgentOS实例中，资格迹可由情景记录 $E$ 保存，工作状态 $h(t)$ 提供即时活动，慢结构 $\Theta$ 接受经过回放和验收的提案，TCE或CCM控制学习资格与预算。这只是工程映射；其他系统可以用局部缓存、事件日志、可微存储或进化搜索实现跨时信用分配。一般结论只有两点：快态必须留下可寻址证据，慢更新必须拥有结果信号与稳定化机制。局部结构节点与全局分布活动通过身份键绑定，若身份丢失，再精确的时间相关也会把信用写入错误对象。

\section{宏观层：系统巩固、回放与结构重组}
\label{sec:macro_cortical_metric_flow_and_consolidation}

跨越一次任务或睡眠周期后，研究对象从单个权重转为多区域协同结构。我们用事件集合 $E=\{e_r\}$ 记录经历，每个事件包含对象身份、时间、上下文、行动、结果和置信度；用 $G_c$ 表示长期关系网络；用分布状态 $h_c$ 表示当前可读出的全局内容。海马—皮层系统的经验事实支持回放与巩固相关，但这不等于海马波包在雕刻测地线。更稳妥的计算描述是：事件索引选择回放样本，回放驱动候选参数更新，系统级门控决定哪些变化被提交。

设 $q_r$ 是事件的回放优先级，依赖惊奇度、任务价值、访问频率、新颖性和覆盖缺口。回放分布为 $P(r)=\operatorname{softmax}(q_r/T_r)$，其中 $T_r$ 是无量纲采样尺度。一次宏观更新写成
\begin{equation}
G_c^{v+1}=\operatorname{Commit}\!\left(G_c^v,\Pi_{\mathcal G}[G_c^v-\eta_c\widehat\nabla_GL_{\rm replay}-\eta_c\lambda_c\nabla_G\Omega],\mathcal A_v\right).
\end{equation}
$v$ 是结构版本，$\Omega$ 约束复杂度和干扰，$\mathcal A_v$ 是接受测试，包括旧任务保持、新任务收益、校准、安全与资源预算。Commit可以拒绝、部分接受或回滚候选结构，因此回放损失下降不能自动推出长期记忆巩固。

尖波涟漪、慢波和睡眠纺锤为生物实现提供候选机制：它们可能协调事件重放、皮层可塑性窗口和跨区域通信。模型应检验三类预测。第一，干预回放顺序是否改变之后的行为和表征；第二，阻断特定节律耦合是否选择性损害跨日保持；第三，控制总活动量后，具有任务相关身份绑定的回放是否比随机回放更有效。若只观察到相关活动而没有干预与读出变化，我们只能提出经验假说，不能断言因果雕刻。

系统巩固也不是把所有情景压成一个平滑全局向量。局部结构网络保留对象、事件和因果依赖，分布表示提供相似检索、泛化与快速读出。二者通过任务相关身份键、绑定算子和版本接口联合：局部节点固定“谁、何时、何关系”，分布状态把跨样本规律展开；读出失败时可回到事件证据，结构冲突时可暂停参数沉淀。海马体与皮层的分工为这种架构提供启发，却不证明工程组件与脑区一一对应。

从费力推理到熟练直觉，可以解释为计算路径缩短：反复成功的事件使中间结果缓存化、结构化或参数化，在线搜索深度、延迟和工作记忆占用下降。但路径缩短并非零能耗，也可能带来僵化。环境改变后，旧快捷路径会产生系统性错误，现实反馈必须重新打开搜索和结构修订。合格的巩固模型同时报告收益、遗忘、迁移、反事实适应和恢复代价，而不是只展示熟悉任务上的速度提升。

工程验证可将回放策略随机分组，保持总样本量与计算预算相同，比较均匀回放、优先回放、按覆盖缺口回放和按价值回放。若某策略只提升训练样本重构而损害新情境迁移，就说明其优化了代理目标。进一步注入标签污染、身份交换和过期事件，可以检验结构提交是否具有来源追踪与回滚能力。AgentOS中的 $E$、$\Theta$ 与外部记忆 $M_e$ 可实现这些接口，但一般生命系统不必采用相同存储介质。

巩固的时间尺度也不能由一次离线阶段预先固定。新事件密度、环境漂移、旧知识脆弱度和可用回放预算会共同改变提交周期。工程上可把候选版本先置于影子读出通道，让它回答同一组新旧任务，再依据置信区间而非单次均值决定晋升；生物实验则可在不同睡眠阶段与清醒休息期追踪同一事件的可解码性。若表征位置改变但任务关系保持，说明发生了再编码，不应误判为遗忘；若关系读出与来源证据同时丢失，才构成更强的结构破坏证据。

\section{种群层：选择、漂变与分布动力学}
\label{sec:population_evolution_information_geometry}

跨世代模型的状态是种群频率而非个体记忆。设离散类型 $a=1,\ldots,m$ 的频率为 $p_a\ge0$ 且 $\sum_a p_a=1$，相对繁殖成功率为 $w_a(p,e)$，环境状态为 $e$。忽略突变和迁移时，复制子方程为
\begin{equation}
 \dot p_a=p_a(w_a(p,e)-\bar w(p,e)),\qquad \bar w=\sum_bp_bw_b.
\end{equation}
它保持概率单纯形，并说明高于平均繁殖成功率的类型频率上升。加入突变矩阵 $Q_{ba}$、迁移 $m_a$ 和有限种群噪声后，漂移项成为 $\sum_bp_bw_bQ_{ba}-p_a\bar w+m_a$，扩散强度通常随有效种群大小 $N_e^{-1/2}$ 缩放。这里的遗传漂变是抽样过程，不是量子隧穿；穿越适应度谷也不要求真实物理势垒的隧穿机制。

在概率单纯形内部，Shahshahani度量或Fisher型局部度量可把复制子动力学表示为特定势函数的自然梯度，但成立条件十分重要：适合度场需能由势函数导出，频率依赖博弈可能存在循环，环境也可能随时间改变。在非势博弈、红皇后动力学或强生态反馈中，平均适合度未必单调增加。Fisher信息描述参数分布的局部可辨识性，不等于突变方向的死亡阻抗；费希尔演化基本定理也不能被扩张为一切进化都沿全局最速路径前进。

如果环境和适合度显式依赖时间，候选目标 $V(p,e,t)$ 的全导数包含 $dV/dt=\nabla_pV^\top\dot p+\nabla_eV^\top\dot e+\partial_tV$。即使第一项沿某种度量下降，环境恶化或目标变化仍可令总量上升。这与个体学习中的变化目标问题相同：局部下降是数学结果，长期适应是经验结论。验证时应报告多样性、平均繁殖成功、灭绝风险、适应速度和环境覆盖，不能用单一适合度替代全部生态后果。

鲍德温效应提供了跨尺度但非拉马克式的案例。个体学习改变行为，行为改变不同基因型在特定环境中的存活与繁殖差异，经过世代选择后，支持更低学习成本或更高可塑性的遗传变异可能增多。可写成 $w_a=\mathbb E[R(a,\ell,e)-C_{\rm learn}(a,\ell,e)]$，其中 $\ell$ 是学习轨迹。个体获得的具体权重不会直接复制到后代；跨代接口是繁殖差异、生态位改变、亲代效应或文化传递。必须区分这些通道，避免把行为反馈夸大为习得结构必然写回基因流形。

反例说明种群适应没有单一方向。小种群中漂变可压倒微弱选择；迁移可能输入不利等位型，也可能提供恢复多样性；短期繁殖优势可能降低长期环境韧性；个体竞争与群体持续性也可能冲突。模型应通过谱系数据、环境操纵和人口结构控制区分这些机制，不能只用频率变化反推选择。选择系数、突变率、迁移率与有效种群大小必须分别估计，并报告不可辨识区间。

频率数据还存在观测层混淆：采样地点改变、世代重叠、隐匿亚群和检测偏差都能制造类似选择的轨迹。因而推断流程应先声明观测模型，再比较中性、选择、迁移和环境切换等候选机制，并用后验预测或留出时间段检查可辨识性。即使某一选择模型拟合最好，也只能在候选集合与采样设计内成立。跨环境复现实验若出现方向相反的选择系数，应优先检查基因型—环境交互，而不是把差异压缩成一个普适的信息几何方向。

种群模型对人工系统只能提供边界清晰的借鉴。多Agent或模型种群可以通过变异、选择和资源重分配更新，但选择单位、复制机制与评价环境必须显式定义。AgentOS的单体记忆更新不等于种群进化，群体策略迭代也不自动具有生物遗传意义。HSF-HD提供状态分布、结构变化和环境反馈的共同记账框架，而不取消个体与种群的层级差异。

% 精确续写点：下一节 sec:hierarchical_coupling_rg_flow
